\ifdefined\gkver\else\def\gkver{student}\fi
\documentclass[\gkver]{gaokaozhenti}
\begin{document}

\chapter{2005年北京春季理综}
%% paperType: 理综物理部分

\begin{enumerate}
\item
%% number: 15
%% paperName: 2005年北京春季理综第 15 题 6 分
%% typeId: 1
%% type: 单选题
%% chapter: 光学
%% point: 光的电磁说
%% method: 无
%% score: 6
%% degree: 850
%% duplicateId: 0
%% body:
有关红、蓝两束单色光，下述说法正确的是\xzanswer{D}
\fourchoices[answer=D]
{在空气中的波长 $\lambda_{\text{红}}<\lambda_{\text{蓝}}$}
{在水中的光速 $v_{\text{红}}<v_{\text{蓝}}$}
{在同一介质中的折射率 $n_{\text{红}}>n_{\text{蓝}}$}
{蓝光光子的能量大于红光光子的能量}
%% answer: D
%% memo:
\memoanswer{
本题原卷及材料中未提供详解，待补充。
}

\item
%% number: 18
%% paperName: 2005年北京春季理综第 18 题 6 分
%% typeId: 1
%% type: 单选题
%% chapter: 电路
%% point: 闭合电路欧姆定律
%% method: 无
%% score: 6
%% degree: 850
%% duplicateId: 0
%% body:
一电池外电路断开时的路端电压为 $3\UV$，接上 $8\UO$ 的负载电阻后路端电压降为 $2.4\UV$，则可以判定电池的电动势 $E$ 和内电阻 $r$ 为\xzanswer{B}
\fourchoices[answer=B]
{$E=2.4\UV$，$r=1\UO$}
{$E=3\UV$，$r=2\UO$}
{$E=2.4\UV$，$r=2\UO$}
{$E=3\UV$，$r=1\UO$}
%% answer: B
%% memo:
\memoanswer{
本题原卷及材料中未提供详解，待补充。
}

\item
%% number: 19
%% paperName: 2005年北京春季理综第 19 题 6 分
%% typeId: 1
%% type: 单选题
%% chapter: 机械波
%% point: 波的图象
%% method: 无
%% score: 6
%% degree: 850
%% duplicateId: 0
%% body:
一列沿 $x$ 轴正方向传播的简谐横波，波速为 $60\Ums$，在 $t=0$ 时波的图像如图所示，则\xzanswer{C}
\onepicture[w=9cm]{figs/19.jpg}
\fourchoices[answer=C]
{此波频率为 $40\UHz$，此时质元 b 的速度为零}
{此波频率为 $40\UHz$，此时质元 b 的速度向着 $y$ 轴负方向}
{此波频率为 $20\UHz$，此时质元 a 的速度向着 $y$ 轴正方向}
{此波频率为 $20\UHz$，此时质元 a 的速度为零}
%% answer: C
%% memo:
\memoanswer{
本题原卷及材料中未提供详解，待补充。
}

\item
%% number: 20
%% paperName: 2005年北京春季理综第 20 题 6 分
%% typeId: 1
%% type: 单选题
%% chapter: 运动和力
%% point: 超重失重
%% method: 无
%% score: 6
%% degree: 650
%% duplicateId: 0
%% body:
如图所示，一个盛水的容器底部有一小孔。静止时用手指堵住小孔不让它漏水，假设容器在下述几种运动过程中始终保持平动，且忽略空气阻力，则\xzanswer{D}
\onepicture[w=4cm]{figs/20.jpg}
\fourchoices[answer=D]
{容器自由下落时，小孔向下漏水}
{将容器竖直向上抛出，容器向上运动时，小孔向下漏水；容器向下运动时，小孔不向下漏水}
{将容器水平抛出，容器在运动中小孔向下漏水}
{将容器斜向上抛出，容器在运动中小孔不向下漏水}
%% answer: D
%% memo:
\memoanswer{
本题原卷及材料中未提供详解，待补充。
}

\item
%% number: 21
%% paperName: 2005年北京春季理综第 21 题 6 分
%% typeId: 1
%% type: 单选题
%% chapter: 电磁感应
%% point: 电磁感应中图象
%% method: 无
%% score: 6
%% degree: 450
%% duplicateId: 0
%% body:
一个边长为 $6\Ucm$ 的正方形金属线框置于匀强磁场中，线框平面与磁场垂直，电阻为 $0.36\UO$。磁感应强度 $B$ 随时间 $t$ 的变化关系如图所示，则线框中感应电流的有效值为\xzanswer{B}
\onepicture[w=8cm]{figs/21.jpg}
\fourchoices[answer=B]
{$\sqrt{2}\times10^{-5}\UA$}
{$\sqrt{6}\times10^{-5}\UA$}
{$\frac{\sqrt{2}}{2}\times10^{-5}\UA$}
{$\frac{3\sqrt{2}}{2}\times10^{-5}\UA$}
%% answer: B
%% memo:
\memoanswer{
本题原卷及材料中未提供详解，待补充。
}

\item
%% number: 22
%% paperName: 2005年北京春季理综第 22 题 12 分
%% typeId: 4
%% type: 实验题
%% chapter: 电路
%% point: 伏安法测电阻
%% method: 无
%% score: 12
%% degree: 450
%% duplicateId: 0
%% body:
\begin{enumerate}
	\item
	图 1 中螺旋测微器的读数为 \tkanswer[2cm]{2.720} $\Umm$。

	\onepicture[w=5.5cm]{figs/22a.jpg}
	\item
	原始的电话机将听筒和话筒串联成一个电路，当自己对着话筒讲话时，会从话筒中听到自己的声音，导致听觉疲劳而影响通话。现代的电话将听筒电路与话筒电路分开，改进的电路原理示意图如图 2 所示，图 2 中线圈Ⅰ与线圈Ⅱ匝数相等，$R_{0}=1.2\UkO$，$R=3.6\UkO$，$R_{x}$ 为可变电阻。当 $R_{x}$ 调到某一值时，从听筒中就听不到话筒传出的声音了，这时 $R_{x}=$ \tkanswer[1.5cm]{1.8} $\UkO$。

	\onepicture[w=6cm]{figs/22b.jpg}
	\item
	在测量电阻的实验中，提供的器材有：$3\UV$ 稳压电源 $E$、滑线变阻器 $R$、电压表 V、电流表 A、待测电阻 $R_{x}$，以及开关 S、导线等。要求：①电流表内接；②调节滑线变阻器可使电压表的示数在 $0\UV\sim3\UV$ 间变化。

	在实验中，有的同学连成图 3、图 4 所示的电路，其中 a，b，c，$\ldots$，k 是表示接线柱的字母。请将图 3、图 4 中接线错误（用导线两端接线柱的字母表示）、引起的后果、改正的方法（改接、撤消或增添），分别填在图下面相应的表格中。

	\onepicture[w=9cm]{figs/22c.jpg}

	\begin{center}
	\begin{tabular}{|c|c|c|}
	\hline
	接线错误 & 引起的后果 & 改正的方法 \\ \hline
	 & & \\ \hline
	 & & \\ \hline
	\end{tabular}
	\end{center}

	\onepicture[w=9cm]{figs/22d.jpg}

	\begin{center}
	\begin{tabular}{|c|c|c|}
	\hline
	接线错误 & 引起的后果 & 改正的方法 \\ \hline
	 & & \\ \hline
	 & & \\ \hline
	\end{tabular}
	\end{center}
\end{enumerate}
%% answer:
\jdanswer*{
\begin{enumerate}
	\item
	$2.720\Umm$。
	\item
	$1.8\UkO$。
	\item
	图 3 对应的表格：
	\begin{center}
	\begin{tabular}{|c|c|c|}
	\hline
	接线错误 & 引起的后果 & 改正的方法 \\ \hline
	ce & 电源可能短路，损坏电源 & 撤销 ce 连线 \\ \hline
	fi，hj & 电流表指针反偏，损坏电流表 & fi$\to$fh，hi$\to$ij \\ \hline
	\end{tabular}
	\end{center}
	图 4 对应的表格：
	\begin{center}
	\begin{tabular}{|c|c|c|}
	\hline
	接线错误 & 引起的后果 & 改正的方法 \\ \hline
	缺 db（或 dg 或 dk）连线 & 电压表的示数无法调到 $0\UV$ & 增添 db（或 dg 或 dk）连线 \\ \hline
	hj & 电流表短路，指针无示数 & 撤销 hi 连线 \\ \hline
	\end{tabular}
	\end{center}
\end{enumerate}
}
%% memo:
\memoanswer{
本题原卷及材料中未提供详解，待补充。
}

\item
%% number: 23
%% paperName: 2005年北京春季理综第 23 题 20 分
%% typeId: 6
%% type: 计算题
%% chapter: 磁场
%% point: 带电粒子在复合场中的运动
%% method: 无
%% score: 20
%% degree: 550
%% duplicateId: 0
%% body:
两块金属 a、b 平行放置，板间存在与匀强电场正交的匀强磁场，假设电场、磁场只存在于两板间的空间区域。一束电子以一定的初速度 $v_{0}$ 从两极板中间，沿垂直于电场、磁场的方向射入场中，无偏转地通过场区，如图所示。

已知板长 $l=10\Ucm$，两板间距 $d=3.0\Ucm$，两板间电势差 $U=150\UV$，$v_{0}=2.0\times10^{7}\Ums$。

（电子所带电荷量的大小与其质量之比 $\frac{e}{m}=1.76\times10^{11}\UCkg$，电子电荷量的大小 $e=1.60\times10^{-19}\UC$）
\begin{enumerate}
	\item
	求磁感应强度 $B$ 的大小；
	\item
	若撤去磁场，求电子穿过电场时偏离入射方向的距离，以及电子通过场区后动能增加多少？
\end{enumerate}
\onepicture[w=9cm, align=right]{figs/23.jpg}
%% answer:
\jdanswer{
\begin{enumerate}
	\item
	$B=2.5\times10^{-4}\UT$。
	\item
	$y_{1}=1.1\times10^{-2}\Um$，$\Delta E_{k}=8.8\times10^{-18}\UJ=55\UeV$。
\end{enumerate}
}
%% memo:
\memoanswer{
（1）电子进入正交的电磁场不发生偏转，则满足 $Bev_{0}=e\frac{U}{d}$，解得 $B=\frac{U}{v_{0}d}=2.5\times10^{-4}\UT$。

（2）设电子通过场区偏转的距离为 $y_{1}$，则 $y_{1}=\frac{1}{2}at^{2}=\frac{1}{2}\cdot\frac{eU}{md}\cdot\frac{l^{2}}{v_{0}^{2}}=1.1\times10^{-2}\Um$；电子通过场区后动能增加 $\Delta E_{k}=eEy_{1}=e\frac{U}{d}y_{1}=8.8\times10^{-18}\UJ=55\UeV$。
}

\item
%% number: 24
%% paperName: 2005年北京春季理综第 24 题 18 分
%% typeId: 6
%% type: 计算题
%% chapter: 运动和力
%% point: 动量守恒定律
%% method: 无
%% score: 18
%% degree: 550
%% duplicateId: 0
%% body:
下雪天，卡车在笔直的高速公路上匀速行驶。司机突然发现前方停着一辆故障车，他将刹车踩到底，车轮被抱死，但卡车仍向前滑行，并撞上故障车，且推着它共同滑行了一段距离 $l$ 后停下。事故发生后，经测量，卡车刹车时与故障车距离为 $L$，撞车后共同滑行的距离 $l=\frac{8}{25}L$。假定两车轮胎与雪地之间的动摩擦因数相同。已知卡车质量 $M$ 为故障车质量 $m$ 的 4 倍。
\begin{enumerate}
	\item
	设卡车与故障车相碰前的速度为 $v_{1}$，两车相撞后的速度变为 $v_{2}$，求 $\frac{v_{1}}{v_{2}}$；
	\item
	卡车司机至少在距故障车多远处采取同样的紧急刹车措施，事故就能免于发生。
\end{enumerate}
%% answer:
\jdanswer{
\begin{enumerate}
	\item
	$\frac{v_{1}}{v_{2}}=\frac{5}{4}$。
	\item
	$L'=\frac{3}{2}L$。
\end{enumerate}
}
%% memo:
\memoanswer{
（1）由碰撞过程动量守恒 $Mv_{1}=(M+m)v_{2}$，又 $M=4m$，则 $\frac{v_{1}}{v_{2}}=\frac{M+m}{M}=\frac{5}{4}$。

（2）设卡车刹车前速度为 $v_{0}$，轮胎与雪地之间的动摩擦因数为 $\mu$。两车相撞前卡车动能变化 $\frac{1}{2}Mv_{0}^{2}-\frac{1}{2}Mv_{1}^{2}=\mu MgL$，碰撞后两车共同向前滑动，动能变化 $\frac{1}{2}(M+m)v_{2}^{2}-0=\mu(M+m)gl$。由以上两式得 $v_{0}^{2}-v_{1}^{2}=2\mu gL$，$v_{2}^{2}=2\mu gl$，又 $l=\frac{8}{25}L$，得 $v_{0}^{2}=3\mu gL$。如果卡车滑到故障车前就停止，由 $\frac{1}{2}Mv_{0}^{2}-0=\mu MgL'$，故 $L'=\frac{3}{2}L$。这意味着卡车司机在距故障车至少 $\frac{3}{2}L$ 处紧急刹车，事故就能够免于发生。
}

\item
%% number: 25
%% paperName: 2005年北京春季理综第 25 题 22 分
%% typeId: 6
%% type: 计算题
%% chapter: 电磁感应
%% point: 电磁感应受力分析
%% method: 无
%% score: 22
%% degree: 450
%% duplicateId: 0
%% body:
近期《科学》中文版的文章介绍了一种新技术——航天飞缆，航天飞缆是用柔性缆索将两个物体连接起来在太空飞行的系统。飞缆系统在太空飞行中能为自身提供电能和拖曳力，它还能清理“太空垃圾”等。从 1967 年至 1999 年 17 次试验中，飞缆系统试验已获得部分成功。该系统的工作原理可用物理学的基本定律来解释。

下图为飞缆系统的简化模型示意图，图中两个物体 P，Q 的质量分别为 $m_{P}$、$m_{Q}$，柔性金属缆索长为 $l$，外有绝缘层，系统在近地轨道作圆周运动，运动过程中 Q 距地面高为 $h$。设缆索总保持指向地心，P 的速度为 $v_{P}$。已知地球半径为 $R$，地面的重力加速度为 $g$。
\begin{enumerate}
	\item
	飞缆系统在地磁场中运动，地磁场在缆索所在处的磁感应强度大小为 $B$，方向垂直于纸面向外。设缆索中无电流，问缆索 P、Q 哪端电势高？此问中可认为缆索各处的速度均近似等于 $v_{P}$，求 P、Q 两端的电势差；
	\item
	设缆索的电阻为 $R_{1}$，如果缆索两端物体 P、Q 通过周围的电离层放电形成电流，相应的电阻为 $R_{2}$，求缆索所受的安培力多大；
	\item
	求缆索对 Q 的拉力 $F_{Q}$。
\end{enumerate}
\onepicture[w=3cm, align=right]{figs/25.jpg}
%% answer:
\jdanswer{
\begin{enumerate}
	\item
	$U_{PQ}=Blv_{P}$，P 点电势高。
	\item
	$F_{A}=\frac{B^{2}l^{2}v_{P}}{R_{1}+R_{2}}$。
	\item
	$F_{Q}=m_{Q}\left[\frac{gR^{2}}{(R+h)^{2}}-\frac{(R+h)v_{P}^{2}}{(R+h+l)^{2}}\right]$。
\end{enumerate}
}
%% memo:
\memoanswer{
（1）缆索的电动势 $E=Blv_{P}$，P、Q 两点电势差 $U_{PQ}=Blv_{P}$，P 点电势高。

（2）缆索电流 $I=\frac{E}{R_{1}+R_{2}}=\frac{Blv_{P}}{R_{1}+R_{2}}$，安培力 $F_{A}=BIl=\frac{B^{2}l^{2}v_{P}}{R_{1}+R_{2}}$。

（3）设 Q 的速度为 $v_{Q}$，Q 受地球引力和缆索拉力 $F_{Q}$ 作用，$G\frac{Mm_{Q}}{(R+h)^{2}}-F_{Q}=m_{Q}\frac{v_{Q}^{2}}{R+h}$；P、Q 角速度相等，$\frac{v_{P}}{v_{Q}}=\frac{R+h+l}{R+h}$；又 $g=\frac{GM}{R^{2}}$，联立解得 $F_{Q}=m_{Q}\left[\frac{gR^{2}}{(R+h)^{2}}-\frac{(R+h)v_{P}^{2}}{(R+h+l)^{2}}\right]$。
}

\end{enumerate}

\end{document}
